\section{Introduction}
\label{sec:intro}

Large language models can write grammatically flawless prose, yet their creative stories are remarkably formulaic: predictable arcs, shallow characters, and repetitive language patterns that readers quickly learn to spot~\cite{chakrabarty2023art}.
This ``formulaic ceiling'' persists even as models scale, suggesting that simply generating more tokens does not produce more creative ones.
A natural question arises: if LLMs struggle with \emph{holistic} story generation, can we decompose the task into manageable components and recombine them to unlock better controllability---or even better quality?

{\bf Why decomposition?}
The \skillmix framework~\cite{allenzhu2024instructskillmix} demonstrated that decomposing instruction-following into discrete skills, creating combinatorial synthetic data from novel skill mixtures, and training on this data significantly improves model capabilities.
Concurrently, decomposition-based story generation systems---DOC~\cite{yang2023doc}, Agents' Room~\cite{huot2024agentsroom}, and DSR~\cite{dsr2025}---have shown that separating planning from writing yields more coherent, relevant narratives.
These findings suggest an intriguing \emph{arbitrage} hypothesis: even if an LLM writes mediocre stories end-to-end, it might produce better stories when given explicit control over narrative dimensions, because the model's limited generation ``bandwidth'' is freed from having to simultaneously discover \emph{and} execute good structural choices.

{\bf The gap.}
Despite the success of both paradigms, no prior work has applied the \skillmix approach of systematic component enumeration, combinatorial synthesis, and controlled evaluation to creative writing.
Existing decomposition systems use ad-hoc structures tailored to specific generation pipelines~\cite{yang2023doc,huot2024agentsroom}; none treat story components as a first-class composable vocabulary analogous to instruction-following skills.

{\bf Our approach.}
We introduce \methodname, a pipeline that: (1) extracts a structured taxonomy of narrative components from existing stories, (2) creates novel specifications by mixing components from different source stories, and (3) generates stories conditioned on these multi-dimensional specifications.
We evaluate on \rocstories~\cite{mostafazadeh2016corpus} using \gptfour for extraction, generation, and evaluation across three conditions: \compcontrolled (full 7-component specification), \titleonly (theme only), and \prompted (theme + setting).

{\bf Key findings.}
Component-controlled generation achieves near-perfect specification adherence (4.96/5.0 overall) but suffers a significant quality penalty: overall quality drops by 0.40 points ($d{=}{-}0.80$, $p{=}0.002$) compared to \titleonly baselines, with the largest degradation in language quality ($d{=}{-}0.83$).
Rather than the hypothesized arbitrage, we observe a \emph{constraint-quality tradeoff}: the model faithfully follows all specifications but produces flatter, less creative prose in the process.
Diversity metrics show modest improvements (26\% larger vocabulary), and a supplementary experiment reveals that the quality penalty stems from the \emph{number} of constraints, not from component incoherence---mixed specifications perform as well as coherent ones.

Our contributions are:
\begin{itemize}[leftmargin=*,itemsep=0pt,topsep=0pt]
    \item We apply the \skillmix paradigm to creative story generation, introducing a pipeline for extracting, combining, and evaluating narrative components as composable units.
    \item We demonstrate near-perfect controllability (4.96/5.0) for multi-dimensional story specifications, establishing that LLMs can reliably follow complex narrative constraints.
    \item We identify a constraint-quality tradeoff in which specification adherence competes with creative expression, refuting the quality arbitrage hypothesis for single-pass generation.
    \item We provide evidence that a two-stage approach---structural specification followed by unconstrained refinement---is needed to realize the benefits of component-based story generation.
\end{itemize}
